{
\parskip=10pt    
\makeabstract
    {
    The Potential of Acetonitrile as a Novel Denaturant
    }
    {
    Sofia Cante\textsuperscript{2}, Beatrice Stefan\textsuperscript{2}, Ryan Tak\textsuperscript{1}, Sheila Jaswal\textsuperscript{1,2}
    }
    {
    \textsuperscript{1} Biochemistry \& Biophysics Program, Amherst College, Amherst, MA \\
\textsuperscript{2} Department of Chemistry, Amherst College, Amherst, MA
    }
    {
    The Central Dogma outlines how the sequence of a protein determines its shape, and how shape determines its function in the body. Proteins move between the most energetically stable conformations depending on their environment. Systematically denaturing a protein can provide insight into the kinetic and thermodynamic factors behind this energetic landscape, supporting research into conformational diseases such as Alzheimer's, Parkinson's, and more. While fluorescence spectroscopy is the traditional method for analyzing protein unfolding, mass spectrometry (MS) is better for analyzing protein dynamics under native conditions in real time. Acetonitrile (AcN) is an ideal medium for protein mass spectrometry due to its volatility, but it must first be proven to be a novel protein denaturant. 
    }
    {
    13.3 {\textmu}M samples of desalted myoglobin (Myb) were prepared with varying AcN concentrations using pH 6.8 Ammonium Acetate buffer (AmAc), incubated for 2 hours at 37{\textdegree}C.  The concentration of Myb was confirmed with UV- Vis. 15 {\textmu}M samples of Trypsinogen (Tgn) were prepared with varying AcN concentrations in $\sim$20 mM, pH 5.1 - 6.85 AmAc buffer, and incubated at 20{\textdegree}C or 25{\textdegree}C for varying amounts of time. pH of buffers was tested using a potassium chloride (KCl) pH electrode.  Fluorescence data measured the intensity and wavelength of emissions from the excited amino acid tryptophan (Trp), which vary as a protein unfolds and reveals more Trp residues. MS sprays the mass to charge ratio (m/z) of protein ions, as charge changes when a protein unfolds. Sprays were collected on a Waters Synapt G2Si Q-TOF Mass Spectrometer.
    }
    {
    Pilot fluorescence of Tgn at 20{\textdegree}C provides a free energy difference of unfolding (\ensuremath{\Delta}Gu) equal to 34.1 - 44.4 kJ/mol, moving toward the values found by traditional denaturations in literature and past AcN denaturations performed by students. Additionally, fluorescence of Tgn over incubation periods ranging from 0 - 119 hours reveals that samples in \ensuremath{\geq}10\% AcN have stable fluorescence data over the first 28 hours of incubation, while samples in low concentrations of AcN (approximately \ensuremath{<}30\%) experience extreme shifts in fluorescence data when incubated for \ensuremath{>}110 hours. Furthermore, Tgn samples in 0\% AcN consistently emit higher average emission wavelengths the longer they are incubated. Mass spectrometry under native conditions and in 0M and 5M AcN displays the presence of native forms of Myb. Mass spectrometry demonstrated that Myb exists in higher charge states at higher AcN concentrations. Simultaneously, the Myb with 9.6M AcN also occupied the apo form (loss of heme). Mass spectrometry of Tgn demonstrates a shift from lower charge states to higher charge states as the protein unfolds, specifically as the most populated charge states move from z = 9 \& 8 to z = 10 \& 9 between 0\% and 45\% AcN. 
    }
    {
    Acetonitrile induces unfolding of both Tgn and Myb. Mass spectrometry detects different states of Myb as it proceeds to further denaturation, and increased relative population in higher charge states as Tgn unfolds. Future mass spectrometry experiments with Tgn in low concentrations of AcN may be able to explain why there is variability in its fluorescence data as incubation time changes. More quantitative studies of the rates at which model proteins fold and unfold are necessary to confirm that mass spectrometry is a valid method for assessing the energetic landscape of myoglobin's protein folding.
    }}

    \newpage
    